% Member A (Shayan). Method -> Memory. Condensed from MEMORY.md; code: agent/memory/m1_summary.py
Our memory (M1, \texttt{agent/memory/m1\_summary.py}) keeps three plain-text items per session.

\paragraph{What is stored.}
(i)~\emph{Key facts}: up to 6 user messages kept word for word as \texttt{(turn, text)}, each clipped
to 300 characters. (ii)~A \emph{summary}: one narrative of at most 150 words, rewritten after every
turn by the same backbone model (same seed) with a dedicated summary prompt. (iii)~\emph{Cited ids}:
document ids the bot cited in an earlier turn \emph{and} had really retrieved or read from memory in
that turn. Retrieved documents, scores, prompts, raw model output and citation flags stay in graph
state and traces only; they can be fetched again, and storing them would turn the bot's guesses into
``facts''.

\paragraph{Write rules.}
A user message becomes a key fact on turn~1, or when it contains an \emph{anchor} that no key fact
holds yet as a whole value. Anchors are matched with regular expressions: section references
(\texttt{section 72}, \texttt{s.~489-F}, \texttt{u/s 302}), the seven Act names and their codes,
dates, years, document ids (\texttt{PPC-379}, \texttt{SHC-140894}) and money amounts
(\emph{loss guard}: a bad rewrite can never lose a date or a section number). Questions that only
point back (``based on what I told you at the start\ldots?'') are not stored. If clipping cuts an
anchor off, it is appended to the line. A 7th fact evicts the oldest one except turn~1.
The model sees the previous summary, the new user message, the bot answer (300 characters) and its
citations. A rewrite is rejected by a \emph{drift guard} if it is empty, contains JSON, exceeds 220
words, or contains a number, section number, document id or Act name absent from the previous
summary, the key facts and the new turn; it is also rejected if the model call fails. The fallback
is the previous summary plus one line \texttt{User (turn N): <message>}, oldest lines dropped past
150 words. Cited ids are added every turn, grounded ones only.

\paragraph{Corrections.}
The summary prompt asks the model to replace corrected values. Independently, when exactly one old and
one new date, section or amount can be identified, the old value is replaced in both the key facts and
the summary. The old value is the stored value following ``not'', ``rather than'' or ``instead of''
(``5 May 2024, not 3 March 2025''), or, if the message says ``correction'', ``I meant'' or ``my
mistake'', the single stored value of that kind. ``Actually'' or ``sorry'' alone never trigger a
replacement. Unclear corrections are stored as a new key fact under a header stating that later lines
override earlier ones.

\paragraph{Read rules.}
The memory text is: a header, the key facts as \texttt{- (turn N) text} (turn~1 first, then oldest to
newest), \texttt{Assistant cited earlier: <ids>} (last 8), then \texttt{Summary:} and the summary.
There is no ranking against the query. Over the 768-token budget, summary lines are dropped first
(oldest first), then key facts other than turn~1; as a last resort the text is cut from the end.
The memory also returns a \emph{retrieval hint} for the fixed BM25 retriever: about 40 topic words
from the key facts only (turn~1, then newest; dates, amounts and chat filler removed), written twice
so the question's own filler words do not outweigh them, plus the document id implied by
``section~N of Act~X'' (e.g.\ \texttt{TPA-72}). The hint is empty when the question names its own
section or document id and does not point back (``this section'', ``my matter''). On the 546 T3/T4
sessions of dev and train, the gold document ranked first in 92\% of cases with the hint (top~5:
99\%), against 81\% with the raw turn-1 text and 0\% with no hint.
